\bookchapter{제26장: 현재 스마트폰 기술 최전선}

\chapterintro{구형 스마트폰의 “사후 처리”에 대해 이야기했으니, 이제 시선을 현재로 돌리자. 2024-2026년은 스마트폰 기술의 집중 폭발기입니다: AI가 클라우드에서 주머니로 들어왔고, 위성 통신이 플래그십에서 중급으로 하향되었으며, 폴더블은 “맛보기”에서 “일상”이 되었고, 배터리 기술도 수년 만에 없던 소재 수준의 돌파구를 맞이. 이 챕터에서는 이러한 진행 중인 기술 최전선을 정리하여 “지금 기기를 교체해야 할지” 그리고 “어떤 기기로 교체할지” 판단하는 데 도움을 드리겠습니다.}

\booksection{AI 스마트폰: 마케팅 허세인가, 진정한 변화인가?}

2024년부터 거의 모든 안드로이드 제조사가 플래그십에 “AI 스마트폰” 깃발을 내걸었고, Apple도 iOS 18에서 Apple Intelligence를 출시. 하지만 어떤 스마트폰이 “AI 스마트폰”이라고 할 수 있을까? 업계는 대략 세 가지 판단 기준을 형성했습니다:

\begin{itemize}
  \item \textbf{단말기 대규모 모델}: 스마트폰 로컬에서 약 7B 파라미터 규모의 언어 모델(예: Qwen, Llama, DeepSeek의 경량 버전)을 실행할 수 있어, 클라우드에 의존하지 않고도 텍스트 요약, 번역, 카피 생성 등 작업을 완료. 판단 방법: 비행 모드에서도 AI 기능을 사용할 수 있다면 단말기 모델임.
  \item \textbf{NPU 컴퓨팅 파워(TOPS)}: NPU(신경망 처리 장치)의 컴퓨팅 파워는 TOPS(초당 조 단위 연산)를 단위로 사용. 2024년 플래그십 칩(예: 스냅드래곤 8 Gen3, 디멘시티 9300)의 NPU 컴퓨팅 파워는 40-70 TOPS 구간이며, 2025-2026년의 차세대 칩은 일반적으로 100 TOPS를 돌파, 이는 단말기 대규모 모델 실행의 하드웨어 문턱.
  \item \textbf{AI Agent 시나리오}: "질의응답"만이 아니라 여러 앱에 걸쳐 작업 체인을 실행——예를 들어 "어제 회의 사진을 정리해서 앨범으로 만들어 장삼에게 공유해 줘"라는 명령에 시스템이 자동으로 앨범, 파일, 위챗 등 여러 앱을 호출하여 완료. 이는 "AI 어시스턴트가 있는 스마트폰"과 "AI 스마트폰"을 구별하는 핵심 표시.
\end{itemize}

\booksection{각 제조사의 AI 전략 대비: 누가 진지하게 하고, 누가 OEM인가?}

각 회사의 AI 노선은 차이가 크므로, 구매 가이드 시 다음에 중점을 둘 것:

\begin{itemize}
  \item \textbf{Apple Apple Intelligence(iOS 18+)}: 시스템 수준 앱 통합 강조——메일 요약, 알림 우선순위, 글쓰기 윤문, 사진관 AI 보정. 단말기 모델이 주를 이루며, 복잡한 작업은 Private Cloud Compute로 클라우드 처리. 장점은 프라이버시 보호와 시스템 원활함, 단점은 중국어 시나리오 적용이 느리고 일부 기능은 2025년에야 점진적으로 지원.
  \item \textbf{Google Gemini Nano}: Pixel 8/9 시리즈 내장, "서클 투 서치", Magic Compose, 녹음 요약 등 지원. Gemini Nano는 진정한 의미의 단말기 소형 모델로, 오픈소스 생태계 친화적이며 안드로이드 AI 스마트폰의 참조 구현.
  \item \textbf{화웨이 판굿 대규모 모델(판구 대규모 모델)}: HarmonyOS의 샤오이(샤오이)에 연결, 단말기 카피 생성, AI 보정, AI 제거(배경 인물 제거), 통화 요약 지원. HarmonyOS NEXT 5.0 이후 시스템 수준 통합이 더욱 강화되어 AI 능력과 시스템의 심층 융합.
  \item \textbf{샤오미 펑박(펑박) 대규모 모델(MiLM)}: HyperOS 내장, AI 사진, AI 요약, AI 화가, 초해상도 지원. 샤오미는 "단말-클라우드 협업" 노선을 채택, 단말기에서 경량 모델 실행, 무거운 작업은 클라우드 처리.
  \item \textbf{OPPO(안디스 대규모 모델)와 vivo(란신 대규모 모델)}: ColorOS 15 및 OriginOS 5에 심층 통합, AI 제거, AI 문서 요약, AI 이미지 인식 질의응답, 자연어 사진 검색("작년에 바닷가에서 찍은 사진 찾아줘") 주력.
\end{itemize}

\textbf{구매 가이드}: 자주 AI로 작업 문서 처리 및 정보 정리가 필요하면 단말기 컴퓨팅 파워가 강하고 AI 기능이 시스템 수준 통합된 기종을 우선 선택; 가끔 AI 보정만 사용한다면 주류 중급 기종으로 충분.

\booksection{위성 통신: 플래그십 “허세”에서 중급 표준으로}

위성 통신은 이 두 해 동안 가장 실질적인 기술 하향 중 하나. 2022년 화웨이 Mate 60과 iPhone 14가 처음 출시했을 때는 플래그십 전용이었지만, 2025년에는 이미 2,000-3,000위안대로 하향.

\begin{itemize}
  \item \textbf{베이두 단문 메시지 중급 기종 하향}: 베이두 위성의 단문 메시지 기능(1,000 한자 이내의 문자 + 위치 전송 가능)이 이미 Redmi, realme, iQOO 등 중급 기종에 하향, 월정액 없이 무료 사용. 야외 트레킹, 자가용 여행 등 시나리오에 적합.
  \item \textbf{천퉁 위성 통화}: 화웨이 Mate 60/70 시리즈, Honor Magic6 시리즈 등은 천퉁 1호 위성 통화 기능 지원, 위성 전화 수발신 및 위성 메시지 전송 가능. 통신사 위성 패키지(약 월 10-200위안) 개통 필요, 무인 지역에 자주 가는 사용자에게 적합.
  \item \textbf{iPhone 위성 메시지(SOS)}: iPhone 14 이후 기종은 위성 SOS 지원, 긴급 서비스 기관에 문자 메시지 및 위치 공유 가능. 2년 무료 서비스 제공, 일부 지역은 갱신 후 사용 가능.
  \item \textbf{실용 조언}: 도시 사용자는 거의 사용하지 않지만, 야외 애호가, 외딴 지역으로 출장하는 비즈니스맨에게는 위성 기능이 있는 기종 선택을 강력 권장——결정적 순간에 생명을 구할 수 있음.
\end{itemize}

\booksection{폴더블 생태계 성숙도: 주력 기기로 사용할 수 있는가?}

폴더블은 “기술 맛보기”에서 “주력 기기 사용 가능”으로 전환. 폴더블의 성숙 여부는 세 가지를 봄: 하드웨어 내구성, 앱 적응도, 상호작용 생태계.

\begin{itemize}
  \item \textbf{하드웨어 내구성}: 힌지 수명이 일반적으로 50만 회 이상(하루 100회 개폐 시 13년 사용 가능), 물방울 힌지가 화면 주름 감소, IPX8 급 방수가 이미 표준. 내부 화면의 UTG 유리+보호막 솔루션으로 화면이 더 이상 "까다롭지 않음".
  \item \textbf{앱 적응}: 주류 앱(위챗, 타오바오, 도우인, B스테이션, Office, WPS)은 이미 대형 화면에 적응되었으나, 약 30\%의 롱테일 앱은 여전히 적응되지 않아 늘어남 또는 검은 테두리로 표시됨. 구매 전에 자주 사용하는 몇 개 앱을 먼저 체험해 볼 것을 권장.
  \item \textbf{현티(서스펜디드, 서스펜디드) 상호작용}: 90°-120°로 접으면 일부 앱은 "상단 화면에 콘텐츠 표시, 하단 화면에서 조작"하는 서스펜디드 모드 지원, 영상 통화, 셀카, 영상 시청에 적합. 카메라 서스펜디드 셀카는 폴더블만의 고유 경험.
  \item \textbf{분할 화면 및 멀티태스킹}: 대형 화면에서 두 개의 앱을 동시에 분할하여 열 수 있으며, 세 개의 앱 소형 창도 공존 가능. HyperOS, ColorOS, OriginOS의 분할 화면 로직은 이미 비교적 성숙, 단 iOS(iPadOS 스타일)의 분할 화면은 iPhone에서 여전히 개방되지 않음.
  \item \textbf{가격}: 횡접이(서적식)는 만위안대에서 5,000-7,000위안대로 내려왔고, 세로 접이(접이식)는 4,000-5,000위안대로 내려옴. 2025년 일부 제조사는 3,000위안대 세로 접이를 출시, 폴더블이 진정으로 대중 시장에 진입.
\end{itemize}

\textbf{구매 가이드}: 무거운 독서, 문서 처리, 영상 사용자라면 횡접이 스크린을 고려할 만함; 휴대성과 사회적 속성을 추구한다면 세로 접이 스크린이 더 적합. 적어도 오프라인 매장에서 30분 만져본 후 결정할 것을 권장.

\booksection{6G와 Wi-Fi 7: 차세대 연결의 속도와 경계}

6G는 아직 본격적으로 착지하지 않았지만, 관련 기술은 이미 진행 중:

\begin{itemize}
  \item \textbf{6G 진전}: 3GPP는 2027년에 6G 표준 제정을 완료하고, 2030년경 상용화할 계획. 6G의 이론적 최고 속도는 1 Tbps(5G의 50배)에 달하며, 지연 시간은 0.1ms까지 낮아져 "통신-감지 통합"(통신+감지)을 지원, 스마트폰이 저정밀 레이더로도 사용 가능. 현재 각 제조사는 여전히 프로토타입 검증 단계에 있으며, 소비자용 제품과는 거리가 멈.
  \item \textbf{Wi-Fi 7(802.11be)}: 2024년부터 플래그십 표준, 이론적 최고 속도 46 Gbps(Wi-Fi 6의 4배), 320MHz 대역폭, 4K-QAM, MLO(다중 링크 동시 전송) 지원. Wi-Fi 7 라우터에서 실제 테스트 시 플래그십 전송 속도 2-5 Gbps, Wi-Fi 6E보다 2-3배 빠름. Wi-Fi 7 스마트폰으로 교체한 경우, 모든 성능을 발휘하려면 Wi-Fi 7 라우터도 함께 구성해야 함.
  \item \textbf{Bluetooth 5.4와 LE Audio}: 2024-2026년 신형 기종은 일반적으로 Bluetooth 5.4 및 LE Audio(저전력 오디오) 지원, 두 개의 이어폰을 동시에 연결해 같은 노래를 들을 수 있으며, Auracast 방송 오디오는 여러 사람이 공공장소에서 같은 오디오 소리를 동기적으로 들을 수 있게 함.
\end{itemize}

\booksection{실리콘-탄소 음극 배터리: 배터리 수명 드디어 맞이한 소재 수준의 돌파}

스마트폰 배터리 수명은 수년간 “하루 한 번 충전”이라는 불만으로 비난받았으나, 2024년부터 실리콘-탄소 음극 배터리 기술의 보급으로 상황이 실질적으로 개선됨.

\begin{itemize}
  \item \textbf{기술 원리}: 전통 흑연 음극의 이론 용량은 약 372 mAh/g이며, 실리콘의 이론 용량은 4200 mAh/g(약 10배)에 달하나, 실리콘은 충전·방전 시 부피 팽창이 300\%에 달해 쉽게 분말화됨. 실리콘-탄소 음극은 나노 실리콘-탄소 복합 구조로 실리콘의 팽창 문제를 허용 가능한 범위로 제어, 배터리 에너지 밀도를 15\%-30\% 향상.
  \item \textbf{실제 성능}: 2024-2025년 실리콘-탄소 음극을 채택한 플래그십(예: 샤오미 14 Ultra, Honor Magic 6 울트라 에디션, vivo X100 Pro 등)은 7-8mm 본체에 5,000-6,000 mAh 배터리를 장착, 이전 세대보다 약 800-1,000 mAh 더 많아 무거운 사용 시 하루 반 정도 사용 가능.
  \item \textbf{백와트 급속 충전 협력}: 실리콘-탄소 음극은 급속 충전과 호환성이 양호하며, 100W/120W 급속 충전으로 17분에 충전 완료가 이미 뉴스가 아님. 2025년 일부 기종은 150W-240W 지원, 10분 내에 80\%까지 충전.
  \item \textbf{구매 가이드}: 무거운 사용자라면 "실리콘-탄소 음극+백와트 급속 충전" 조합을 채택한 2024년 이후 기종을 우선 선택, 배터리 수명 불안이 크게 완화될 것.
\end{itemize}

\booksection{본 챕터 요약: 한 문장으로 이러한 최전선 기억하기}

\begin{itemize}
  \item AI 스마트폰 세 가지 기준: 단말기 대규모 모델, NPU 컴퓨팅 파워 TOPS, AI Agent 다중 앱 시나리오.
  \item 위성 통신: 베이두 단문 메시지 중급 하향, 천퉁 통화는 여전히 플래그십.
  \item 폴더블: 하드웨어 성숙, 앱 적응 70\%, 서스펜디드/분할 화면은 고유 경험.
  \item 6G는 아직 진행 중, Wi-Fi 7은 이미 착지, Bluetooth LE Audio 보급 시작.
  \item 실리콘-탄소 음극+백와트 급속 충전: 배터리 수명과 충전의 이중 돌파.
\end{itemize}

\begin{tcbpractice}{도전: 다음 스마트폰을 위한 "기술 요구 목록" 작성}
종이 한 장을 꺼내 “필수/가산/불필요” 세 단계로 AI, 위성, 폴더블, Wi-Fi 7, 급속 충전, 배터리 용량에 대한 요구를 적어보세요. 이렇게 하면 다음에 오프라인 매장에 갔을 때, 가이드의 “판매 포인트”가 진짜 요구인지 판매 권유인지 즉시 판단할 수 있습니다. 예:

\begin{itemize}
  \item 필수: 5,000mAh 이상 배터리, 80W 이상 급속 충전, NFC, IP68 방수.
  \item 가산: 단말기 AI 요약, Wi-Fi 7, 베이두 단문 메시지.
  \item 불필요: 폴더블(나는 직사각형에 익숙), 위성 통화(멀리 안 감), 240W 극한 급속 충전.
\end{itemize}

\end{tcbpractice}

\begin{tcbsummary}{본 챕터 요약}
\begin{itemize}
  \item AI 스마트폰 세 가지 판단 기준: 단말기 대규모 모델, NPU TOPS, AI Agent 시나리오
  \item 각 제조사 AI 전략: Apple 프라이버시 통합 중시, Google 단말기 오픈소스, 화웨이/샤오미/OPPO/vivo 각자 장점
  \item 위성 통신: 베이두 단문 메시지 중급기 하향, 천퉁 통화는 여전히 플래그십 전용
  \item 폴더블: 하드웨어 이미 성숙, 앱 적응 약 70\%, 서스펜디드와 분할 화면이 고유 경험
  \item 6G 표준 2027년 제정, Wi-Fi 7 이미 보급, Bluetooth LE Audio 착지 시작
  \item 실리콘-탄소 음극 배터리가 5,000-6,000mAh를 플래그십 표준으로 만들었으며, 배터리 수명 불안이 크게 완화
\end{itemize}

\end{tcbsummary}
